\section{The model}\label{sec:model}

This section defines QAM without reference to any tool. The \LaTeX{} package, the Lean checker and the database of later sections
are implementations of this model and apply the same rules.

\subsection{Records, frames and maps}

\begin{definition}[Record]
A \emph{record} is a finite sequence of entries $e_1,\dots,e_n$. It is extended only at the end (\emph{append-only}); an existing
entry is never changed. A correction is a new entry that names the entry it replaces.
\end{definition}

A \emph{frame} is the name of a frame of description: an observer, a convention, a choice of units, a model assumption, a sector of a
theory. Frames have no internal structure in the model. QAM knows only their names and the maps declared between them. A \emph{map}
$g\colon A\to B$ is a declared name with a source frame and a target frame; two maps can be declared to compose to a third. One
frame, written $*$, is reserved for convention-free statements such as a pure number or a theorem of mathematics; it is compatible
with every frame. If an application has several independent convention axes (say, an observer and a choice of units), each
combination gets a frame of its own, with maps between them.

In the example of Section~\ref{sec:example} there are two frames, Lab and Boosted, and one map, the boost along $z$, which we
call $\mathit{lorentz}_z\colon\mathrm{Lab}\to\mathrm{Boosted}$.

\subsection{Entries}

\begin{definition}[Entry]\label{def:entry}
A statement entry is a tuple
\[
e=(\mathit{kind},\ \mathit{key},\ \mathit{frame},\ \mathit{status},\ \mathit{inputs},\ \mathit{premises},\ \mathit{statement})
\]
where
\begin{itemize}[nosep,leftmargin=1.5em]
\item $\mathit{kind}$ is one of \texttt{input}, \texttt{claim}, \texttt{thesis}, \texttt{derive}, \texttt{convert};
\item $\mathit{key}$ is a name that is unique in the record;
\item $\mathit{frame}$ is a declared frame;
\item $\mathit{status}$ is taken from the totally ordered set
\[
\st{open}<\st{reading}<\st{conditional}<\st{computed}<\st{proved};
\]
\item $\mathit{inputs}$ is the set of input keys the entry rests on;
\item $\mathit{premises}$ is the list of earlier entries it was obtained from (for \texttt{convert}, together with the map used);
\item $\mathit{statement}$ is an optional text.
\end{itemize}
\end{definition}

The five kinds have fixed meanings. An \emph{input} is a value put in from outside: a measurement, a value from the literature, a
fitted parameter. A \emph{claim} is a statement without premises in the record, such as a cited theorem or an assumption. A
\emph{thesis} is a claim whose wording is part of its identity; a piece of work typically starts from one. A \emph{derivation} obtains
a statement from premises in the same frame, and a \emph{conversion} carries a statement into another frame through a declared map.
Two further entries carry no status: a \emph{comparison} of two entries, and a \emph{declaration of independence}, which asserts
that an entry rests on no input.

The status levels mean, from strongest to weakest: \st{proved}, established by proof; \st{computed}, supported by a calculation but
not proved; \st{conditional}, valid under named conditions that are not discharged; \st{reading}, an interpretation or assignment,
for instance the identification of a mathematical object with a physical quantity; \st{open}, not established. The author assigns
the status. The model does not check whether something declared \st{proved} has in fact been proved. It checks only whether a
conclusion is declared stronger than what it rests on.

In the example, the preparation~\eqref{eq:prep} together with the boost speed is a thesis (a claim with fixed wording) in Lab with status \st{conditional}: it is an
assumption about the experiment. The measured interval for $t$ is an input in Lab. The narrowed prediction~\eqref{eq:interval} is
obtained by first deriving, in Lab, a statement from the preparation and the measurement, and then converting it to Boosted through
$\mathit{lorentz}_z$.

\subsection{Rules}

\begin{definition}[Well-formed record]\label{def:rules}
A record is \emph{well formed} if every entry $e_i$ satisfies:
\begin{enumerate}[nosep,label=(R\arabic*),leftmargin=2.5em]
\item \emph{Declaration before use.} Every key, frame and map that $e_i$ mentions is declared in $e_1,\dots,e_{i-1}$, and no key is
declared twice.
\item \emph{No promotion.} If $e_i$ is a derivation or conversion with premises $p_1,\dots,p_m$, then
$\mathit{status}(e_i)\le\min_j\mathit{status}(p_j)$.
\item \emph{Inputs are preserved.} An input has $\mathit{inputs}=\{\mathit{key}\}$; a claim and a thesis have
$\mathit{inputs}=\emptyset$; for a derivation or conversion, $\mathit{inputs}(e_i)=\bigcup_j\mathit{inputs}(p_j)$. A declaration of
independence for $e$ is admissible only if $\mathit{inputs}(e)=\emptyset$.
\item \emph{No frame mixing.} The premises of a derivation and the two sides of a comparison lie in the same frame (or in $*$). A
conversion $\mathrm{convert}(g,p)$ with $g\colon A\to B$ requires $\mathit{frame}(p)=A$ and has frame $B$.
\end{enumerate}
\end{definition}

The three substantive rules answer the three silent transitions: (R2) prevents promotion, (R3) prevents an input from passing as a
result, and (R4) prevents frame mixing. (R1) makes it possible to check a record in a single pass from top to bottom.

In the example, (R2) forbids declaring the narrowed prediction \st{proved}, because it rests on a \st{conditional} preparation. (R3)
makes the measured input appear in the input set of the prediction, so that a declaration that the prediction is independent of
inputs is rejected. (R4) forbids comparing the Lab register~\eqref{eq:rho-lab} directly with the Boosted register~\eqref{eq:rho-boosted}.
Section~\ref{sec:latex} shows the actual error messages.

\subsection{Open conditions}\label{sec:open}

The stored fields determine, for every statement, what it ultimately rests on. We use the following terminology.

\begin{definition}[Open conditions]\label{def:open}
The \emph{open conditions} of an entry $e$ are the entries reachable from $e$ along premises whose status is below \st{proved},
together with $\mathit{inputs}(e)$.
\end{definition}

This is a reading of the stored structure, not an additional field. Section~\ref{sec:properties} proves that open conditions are
\emph{inherited}: a statement that rests on anything below \st{proved} is not \st{proved} itself, and every input remains visible however
many steps are taken.
An open condition therefore cannot disappear because further conclusions are built on top of it. Open conditions are objects of the
dependency graph, not remarks in the prose.

\subsection{Identity through a hash chain}

\begin{definition}[Entry hash]\label{def:hash}
The hash of an entry is
\[
h(e)=\mathrm{MD5}\bigl(\mathit{kind};\mathit{key};\mathit{frame};\mathit{status};\mathit{inputs};\mathit{sources};\mathit{statement}\bigr),
\]
where $\mathit{sources}$ is the list of the hashes of the premises of $e$. The hash of the record is the MD5 of the sequence of all
entry hashes.
\end{definition}

Because the hashes of the premises enter the hash of a conclusion, every conclusion names the exact versions of what it rests on. If
one word of a thesis changes, every hash that depends on it changes; independent entries keep theirs. Equal hashes mean equal
entries, not true ones. MD5 serves here to detect accidental change, not deliberate forgery; a release adds SHA-256 sums over the
source, the typeset document and the exported record.

\subsection{What the model leaves out}

The model has no semantics for statements. It does not know what a formula means, whether a calculation is right, or whether a frame
was chosen sensibly. Maps carry no status: if the justification of a map between two frames is itself a condition, it must be recorded
as an entry of its own. A claim declared \st{proved} without premises is accepted as such. A fitted value that is declared again under
a new name as a claim is not recognised as an input. These limits are deliberate. They keep the model small enough that its rules can
be checked completely and mechanically, by a compiler as well as by a proof assistant.
